% ========================================================
% BÀI TẬP TOÁN 5 - KẾT NỐI TRI THỨC VỚI CUỘC SỐNG
% Bài 43: Thực hành và trải nghiệm sử dụng máy tính cầm tay
% Định dạng: BTTeX Part III - Trắc nghiệm điền số ngắn
% ========================================================

% Câu 1
\begin{ex}
  Sử dụng máy tính cầm tay tính giá trị của biểu thức: $16 \times 20$.
  \shortans{320}
  \loigiai{
    Bấm phím [16] [x] [20] [=] ta được kết quả: 320.
  }
\end{ex}

% Câu 2
\begin{ex}
  Sử dụng máy tính cầm tay tính giá trị của biểu thức: $17 \times 20$.
  \shortans{340}
  \loigiai{
    Bấm phím [17] [x] [20] [=] ta được kết quả: 340.
  }
\end{ex}

% Câu 3
\begin{ex}
  Sử dụng máy tính cầm tay tính giá trị của biểu thức: $18 \times 20$.
  \shortans{360}
  \loigiai{
    Bấm phím [18] [x] [20] [=] ta được kết quả: 360.
  }
\end{ex}

% Câu 4
\begin{ex}
  Sử dụng máy tính cầm tay tính giá trị của biểu thức: $19 \times 20$.
  \shortans{380}
  \loigiai{
    Bấm phím [19] [x] [20] [=] ta được kết quả: 380.
  }
\end{ex}

% Câu 5
\begin{ex}
  Sử dụng máy tính cầm tay tính giá trị của biểu thức: $20 \times 20$.
  \shortans{400}
  \loigiai{
    Bấm phím [20] [x] [20] [=] ta được kết quả: 400.
  }
\end{ex}

% Câu 6
\begin{ex}
  Sử dụng máy tính cầm tay tính giá trị của biểu thức: $21 \times 20$.
  \shortans{420}
  \loigiai{
    Bấm phím [21] [x] [20] [=] ta được kết quả: 420.
  }
\end{ex}

% Câu 7
\begin{ex}
  Sử dụng máy tính cầm tay tính giá trị của biểu thức: $22 \times 20$.
  \shortans{440}
  \loigiai{
    Bấm phím [22] [x] [20] [=] ta được kết quả: 440.
  }
\end{ex}

% Câu 8
\begin{ex}
  Sử dụng máy tính cầm tay tính giá trị của biểu thức: $23 \times 20$.
  \shortans{460}
  \loigiai{
    Bấm phím [23] [x] [20] [=] ta được kết quả: 460.
  }
\end{ex}

% Câu 9
\begin{ex}
  Sử dụng máy tính cầm tay tính giá trị của biểu thức: $24 \times 20$.
  \shortans{480}
  \loigiai{
    Bấm phím [24] [x] [20] [=] ta được kết quả: 480.
  }
\end{ex}

% Câu 10
\begin{ex}
  Sử dụng máy tính cầm tay tính giá trị của biểu thức: $25 \times 20$.
  \shortans{500}
  \loigiai{
    Bấm phím [25] [x] [20] [=] ta được kết quả: 500.
  }
\end{ex}

% Câu 11
\begin{ex}
  Sử dụng máy tính cầm tay tính giá trị của biểu thức: $26 \times 20$.
  \shortans{520}
  \loigiai{
    Bấm phím [26] [x] [20] [=] ta được kết quả: 520.
  }
\end{ex}

% Câu 12
\begin{ex}
  Sử dụng máy tính cầm tay tính giá trị của biểu thức: $27 \times 20$.
  \shortans{540}
  \loigiai{
    Bấm phím [27] [x] [20] [=] ta được kết quả: 540.
  }
\end{ex}

% Câu 13
\begin{ex}
  Sử dụng máy tính cầm tay tính giá trị của biểu thức: $28 \times 20$.
  \shortans{560}
  \loigiai{
    Bấm phím [28] [x] [20] [=] ta được kết quả: 560.
  }
\end{ex}

% Câu 14
\begin{ex}
  Sử dụng máy tính cầm tay tính giá trị của biểu thức: $29 \times 20$.
  \shortans{580}
  \loigiai{
    Bấm phím [29] [x] [20] [=] ta được kết quả: 580.
  }
\end{ex}

% Câu 15
\begin{ex}
  Sử dụng máy tính cầm tay tính giá trị của biểu thức: $30 \times 20$.
  \shortans{600}
  \loigiai{
    Bấm phím [30] [x] [20] [=] ta được kết quả: 600.
  }
\end{ex}

% Câu 16
\begin{ex}
  Sử dụng máy tính cầm tay tính giá trị của biểu thức: $31 \times 20$.
  \shortans{620}
  \loigiai{
    Bấm phím [31] [x] [20] [=] ta được kết quả: 620.
  }
\end{ex}

% Câu 17
\begin{ex}
  Sử dụng máy tính cầm tay tính giá trị của biểu thức: $32 \times 20$.
  \shortans{640}
  \loigiai{
    Bấm phím [32] [x] [20] [=] ta được kết quả: 640.
  }
\end{ex}

% Câu 18
\begin{ex}
  Sử dụng máy tính cầm tay tính giá trị của biểu thức: $33 \times 20$.
  \shortans{660}
  \loigiai{
    Bấm phím [33] [x] [20] [=] ta được kết quả: 660.
  }
\end{ex}

% Câu 19
\begin{ex}
  Sử dụng máy tính cầm tay tính giá trị của biểu thức: $34 \times 20$.
  \shortans{680}
  \loigiai{
    Bấm phím [34] [x] [20] [=] ta được kết quả: 680.
  }
\end{ex}

% Câu 20
\begin{ex}
  Sử dụng máy tính cầm tay tính giá trị của biểu thức: $35 \times 20$.
  \shortans{700}
  \loigiai{
    Bấm phím [35] [x] [20] [=] ta được kết quả: 700.
  }
\end{ex}

% Câu 21
\begin{ex}
  Sử dụng máy tính cầm tay tính giá trị của biểu thức: $36 \times 20$.
  \shortans{720}
  \loigiai{
    Bấm phím [36] [x] [20] [=] ta được kết quả: 720.
  }
\end{ex}

% Câu 22
\begin{ex}
  Sử dụng máy tính cầm tay tính giá trị của biểu thức: $37 \times 20$.
  \shortans{740}
  \loigiai{
    Bấm phím [37] [x] [20] [=] ta được kết quả: 740.
  }
\end{ex}

% Câu 23
\begin{ex}
  Sử dụng máy tính cầm tay tính giá trị của biểu thức: $38 \times 20$.
  \shortans{760}
  \loigiai{
    Bấm phím [38] [x] [20] [=] ta được kết quả: 760.
  }
\end{ex}

% Câu 24
\begin{ex}
  Sử dụng máy tính cầm tay tính giá trị của biểu thức: $39 \times 20$.
  \shortans{780}
  \loigiai{
    Bấm phím [39] [x] [20] [=] ta được kết quả: 780.
  }
\end{ex}

% Câu 25
\begin{ex}
  Sử dụng máy tính cầm tay tính giá trị của biểu thức: $40 \times 20$.
  \shortans{800}
  \loigiai{
    Bấm phím [40] [x] [20] [=] ta được kết quả: 800.
  }
\end{ex}
